% from ega3-7-fr.tex
\oldpage[III]{212}
on ait $\mathrm R^if_*(\mathscr F(n))=0$ pour tout $i>0$ (2.2.1) ; la conclusion
résulte donc de (7.9.10).
\end{proof}

Le critère de platitude suivant sera important dans la théorie des « modules »
des faisceaux cohérents du chap. V :

\begin{proposition}[7.9.14]
\label{III.7.9.14.fr}
Soient $Y$ un préschéma noethérien, $f:X\to Y$ un morphisme projectif,
$\mathscr L$ un $\mathscr O_X$-Module inversible ample pour $f$, et posons
$\mathscr F(n)=\mathscr F\otimes\mathscr L^{\otimes n}$ pour tout
$\mathscr O_X$-Module $\mathscr F$ et tout $n\in\mathbf Z$. Pour qu'un
$\mathscr O_X$-Module cohérent $\mathscr F$ soit $Y$-plat, il faut et il suffit
qu'il existe un entier $n_0$ tel que, pour tout $n\geq n_0$,
$f_*(\mathscr F(n))$ soit un $\mathscr O_Y$-Module localement libre.
\end{proposition}

\begin{proof}
La nécessité de la condition se démontre comme dans (7.9.13) (le résultat de
(2.2.1) s'appliquant à un faisceau ample $\mathscr L$, puisque $f$ est
projectif). Pour démontrer la réciproque, on peut se borner au cas où $Y$ est
affine d'anneau $A$ ; en vertu de l'hypothèse et de (2.2.2, (i)), les
$A$-modules $\Gamma(X,\mathscr F(n))$ sont de type fini et projectifs (Bourbaki,
\emph{Alg. comm.}, chap. II, \S~5, n\textsuperscript{o}~2, th. 1). Soit $S$
l'anneau gradué
\[
  S=\bigoplus_{n\geq0}\Gamma(X,\mathscr L^{\otimes n}) ;
\]
on sait que $X$ s'identifie canoniquement à $\operatorname{Proj}(S)$ (II,
4.5.2, (b) et 5.4.4). Soit
\[
  M=\bigoplus_{n\geq n_0}\Gamma(X,\mathscr F(n)) ;
\]
remplaçant au besoin $\mathscr L$ par une puissance $\mathscr L^{\otimes d}$,
on peut supposer que $S$ est engendré par un nombre fini d'éléments de degré 1
(2.3.5.1), et il résulte alors de (II, 2.7.5 et 2.7.2) que $\mathscr F$
s'identifie à $\mathscr P\!roj_0(M)$. Pour tout élément homogène $g\in S$ de
degré $>0$, on a donc $\Gamma(X_g,\mathscr F)=M_{(g)}$ ; or, $M$, somme directe
de $A$-modules projectifs, est un $A$-module plat, donc il en est de même de
$M_g$ ($0_{\mathrm I}$, 6.3.2), et par suite aussi de $M_{(g)}$, qui est un
facteur direct de $M_g$ ($0_{\mathrm I}$, 6.1.2). On en conclut (I, 4.14.5) que
$\mathscr F$ est $Y$-plat en tout point de $X_g$, et comme les $X_g$ recouvrent
$X$, la proposition est démontrée.
\end{proof}

\bigskip
\noindent\emph{(A suivre.)}

\clearpage
